% Автоматаар үүсгэсэн — эх файл: lab06/README.md
\chapter{Лаб 6 — Ирмэгийн хиймэл оюуны байршуулалт ба гүйцэтгэлийн хэмжилт}

\begingroup\scriptsize\begin{longtable}[]{@{}ll@{}}
\toprule
\endhead
\bottomrule
\endlastfoot
\textbf{7 хоног} & XII \\
\textbf{Хугацаа} & 4 цаг \\
\textbf{Суралцахуйн үр дүн} & ҮД5 (шинжлэх), ҮД7 (хэмжих) \\
\textbf{Үнэлгээ} & Бичгийн тайлан, 10 оноо \\
\end{longtable}\endgroup

\hypertarget{ux437ux43eux440ux438ux43bux433ux43e}{%
\section{1. Зорилго}\label{ux437ux43eux440ux438ux43bux433ux43e}}

⚠ \textbf{Энэ лаборатори загвар СУРГАХ тухай биш.} Сургалтыг X--XI долоо хоногийн бие даалтаар Edge Impulse Studio дээр аль хэдийн хийсэн байх ёстой.

Энд бид \textbf{байршуулалтын өртгийг} хэмжинэ:

\begin{verbatim}
  нарийвчлал  ←→  саатал  ←→  санах ой  ←→  эрчим хүч
\end{verbatim}

Загвар 97\% нарийвчлалтай байж болно --- гэвч дүгнэлт нь 800 мс авдаг бол 100 Гц мэдрэгчид \textbf{ашиггүй}. Ирмэгийн ML-ийн жинхэнэ ажил бол энэ дөрвөн хэмжигдэхүүний хооронд ухаалаг солилцоо хийх явдал.

\hypertarget{ux445ux430ux440ux438ux443ux43bux430ux445-ux451ux441ux442ux43eux439-ux430ux441ux443ux443ux43bux442ux443ux443ux434}{%
\subsection{Хариулах ёстой асуултууд}\label{ux445ux430ux440ux438ux443ux43bux430ux445-ux451ux441ux442ux43eux439-ux430ux441ux443ux443ux43bux442ux443ux443ux434}}

\begin{enumerate}
\def\labelenumi{\arabic{enumi}.}
\tightlist
\item
  int8 квантчилал саатлыг хэдэн дахин багасгав? Нарийвчлалыг хэр алдав?
\item
  Дүгнэлтийг ирмэг дээр хийх нь үүлэн рүү илгээхээс хэзээ \textbf{дээр} вэ?
\item
  4 цөмийг бүгдийг ашиглах нь үргэлж дээр үү?
\end{enumerate}

\hypertarget{ux443ux440ux44cux434ux447ux438ux43bux441ux430ux43d-ux43dux4e9ux445ux446ux4e9ux43b-ux431ux438ux435-ux434ux430ux430ux43bux442ux430ux430ux440-ux443ux440ux44cux434ux447ux438ux43bux430ux43d-ux445ux438ux439ux43dux44d}{%
\section{2. Урьдчилсан нөхцөл --- БИЕ ДААЛТААР УРЬДЧИЛАН ХИЙНЭ}\label{ux443ux440ux44cux434ux447ux438ux43bux441ux430ux43d-ux43dux4e9ux445ux446ux4e9ux43b-ux431ux438ux435-ux434ux430ux430ux43bux442ux430ux430ux440-ux443ux440ux44cux434ux447ux438ux43bux430ux43d-ux445ux438ux439ux43dux44d}}

Лабораторийн цагаар загвар сургах хугацаа \textbf{байхгүй}. Дараах зүйлс бэлэн байх ёстой:

\hypertarget{x-ux434ux43eux43bux43eux43e-ux445ux43eux43dux43eux433-ux4e9ux433ux4e9ux433ux434ux4e9ux43b-ux446ux443ux433ux43bux443ux443ux43bux430ux445}{%
\subsection{X долоо хоног --- өгөгдөл цуглуулах}\label{x-ux434ux43eux43bux43eux43e-ux445ux43eux43dux43eux433-ux4e9ux433ux4e9ux433ux434ux4e9ux43b-ux446ux443ux433ux43bux443ux443ux43bux430ux445}}

\begin{Shaded}
\begin{Highlighting}[]
\CommentTok{\# 1) Хэвийн ажиллагааны өгөгдөл}
\ExtensionTok{python3}\NormalTok{ tools/sim\_device.py }\AttributeTok{{-}{-}host} \VariableTok{$H} \AttributeTok{{-}{-}devices}\NormalTok{ 1 }\AttributeTok{{-}{-}interval}\NormalTok{ 0.1 }\KeywordTok{\&}
\ExtensionTok{python3}\NormalTok{ lab06/collect\_data.py }\AttributeTok{{-}{-}host} \VariableTok{$H} \AttributeTok{{-}{-}label}\NormalTok{ normal }\AttributeTok{{-}{-}seconds}\NormalTok{ 300}

\CommentTok{\# 2) Гажилтай өгөгдөл}
\ExtensionTok{python3}\NormalTok{ tools/sim\_device.py }\AttributeTok{{-}{-}host} \VariableTok{$H} \AttributeTok{{-}{-}devices}\NormalTok{ 1 }\AttributeTok{{-}{-}interval}\NormalTok{ 0.1 }\AttributeTok{{-}{-}anomaly{-}rate}\NormalTok{ 1.0 }\KeywordTok{\&}
\ExtensionTok{python3}\NormalTok{ lab06/collect\_data.py }\AttributeTok{{-}{-}host} \VariableTok{$H} \AttributeTok{{-}{-}label}\NormalTok{ anomaly }\AttributeTok{{-}{-}seconds}\NormalTok{ 120}
\end{Highlighting}
\end{Shaded}

\hypertarget{xi-ux434ux43eux43bux43eux43e-ux445ux43eux43dux43eux433-edge-impulse-ux434ux44dux44dux440-ux441ux443ux440ux433ux430ux445}{%
\subsection{XI долоо хоног --- Edge Impulse дээр сургах}\label{xi-ux434ux43eux43bux43eux43e-ux445ux43eux43dux43eux433-edge-impulse-ux434ux44dux44dux440-ux441ux443ux440ux433ux430ux445}}

\begin{enumerate}
\def\labelenumi{\arabic{enumi}.}
\tightlist
\item
  https://studio.edgeimpulse.com --- төслөө үүсгэнэ
\item
  \textbf{Data acquisition → Upload data} → \texttt{lab06/data/*.csv}, шошготой нь
\item
  \textbf{Impulse design}: Time series data (цонх 1000 мс, алхам 500 мс) → Spectral Analysis → \textbf{Anomaly Detection (K-means)} эсвэл Classification
\item
  \textbf{EON Tuner} ажиллуулж хамгийн жижиг, хамгийн хурдан хувилбарыг олно
\item
  \textbf{Deployment} → дараах \textbf{гурвуулангийг} татаж авна:

  \begin{itemize}
  \tightlist
  \item
    \texttt{TensorFlow\ Lite\ (float32)} → \texttt{lab06/models/model\_float32.tflite}
  \item
    \texttt{TensorFlow\ Lite\ (int8\ quantized)} → \texttt{lab06/models/model\_int8.tflite}
  \item
    \texttt{Linux\ (AARCH64)} → \texttt{lab06/models/anomaly.eim}
  \end{itemize}
\item
  Studio-гийн \textbf{Model testing} хуудсаас нарийвчлалыг тэмдэглэнэ (Хүснэгт 1-д хэрэгтэй)
\end{enumerate}

\hypertarget{pi-ux434ux44dux44dux440-ux43eux440ux447ux438ux43d-ux431ux44dux43bux434ux44dux445}{%
\subsection{Pi дээр орчин бэлдэх}\label{pi-ux434ux44dux44dux440-ux43eux440ux447ux438ux43d-ux431ux44dux43bux434ux44dux445}}

\begin{Shaded}
\begin{Highlighting}[]
\ExtensionTok{pip3}\NormalTok{ install }\AttributeTok{{-}{-}break{-}system{-}packages}\NormalTok{ numpy ai{-}edge{-}litert edge\_impulse\_linux}
\ExtensionTok{python3} \AttributeTok{{-}c} \StringTok{"from ai\_edge\_litert.interpreter import Interpreter; print(\textquotesingle{}OK\textquotesingle{})"}
\FunctionTok{chmod}\NormalTok{ +x lab06/models/anomaly.eim}
\end{Highlighting}
\end{Shaded}

\begin{quote}
\textbf{Загвар бэлэн болоогүй бол} лабораторид \texttt{-\/-backend\ synthetic} ашиглаж бүх аргачлалыг гүйцэтгэж болно. Оноо хасагдана (3 оноо), гэхдээ лабораторийг алдахгүй.
\end{quote}

\hypertarget{ux43eux43dux43eux43bux44bux43d-ux441ux430ux43dux443ux443ux43bux433ux430}{%
\section{3. Онолын сануулга}\label{ux43eux43dux43eux43bux44bux43d-ux441ux430ux43dux443ux443ux43bux433ux430}}

\textbf{Квантчилал.} float32 жинг int8 болгоно: 4 дахин бага санах ой, бүхэл тоон арифметик тул CPU дээр 2--4 дахин хурдан. Үнэ нь нарийвчлалын бага зэргийн алдагдал (ихэвчлэн \textless{} 1\%).

\textbf{Халаалт (warmup) яагаад чухал вэ.} Эхний хэдэн дүгнэлт үргэлж удаан: санах ойн хуваарилалт, кэш дүүрэх, зарим backend-д JIT. Халаалтын дүгнэлтийг хэмжилтэд оруулбал \textbf{p99 гуйвуулагдана}. \texttt{benchmark\_inference.py} анхдагчаар 30 удаа халаана.

\textbf{p50 биш p99.} Бодит цагийн систем дунджаар биш, \textbf{хамгийн муу тохиолдлоор} төлөвлөгддөг. 100 Гц мэдрэгчид p99 \textless{} 10 мс байх ёстой.

\textbf{Ирмэг vs үүл --- хэзээ аль нь вэ.}

\begingroup\scriptsize\begin{longtable}[]{@{}lll@{}}
\toprule
Шалгуур & Ирмэг дээр & Үүлэн дээр \\
\midrule
\endhead
\bottomrule
\endlastfoot
Саатал & 1--50 мс & 50--500 мс (сүлжээнээс хамаарна) \\
Сүлжээгүй үед & ажиллана & ажиллахгүй \\
Зурвасын өргөн & зөвхөн үр дүн & бүх түүхий өгөгдөл \\
Нууцлал & өгөгдөл гарахгүй & өгөгдөл гарна \\
Загвар шинэчлэх & OTA шаардана (Лаб 2!) & шууд \\
Тооцооллын хүч & хязгаартай & бараг хязгааргүй \\
\end{longtable}\endgroup

\hypertarget{ux430ux43bux445ux43cux443ux443ux434}{%
\section{4. Алхмууд}\label{ux430ux43bux445ux43cux443ux443ux434}}

\hypertarget{ux430ux43bux445ux430ux43c-1-ux441ux443ux443ux440ux44c-ux431ux430-ux430ux440ux433ux430ux447ux43bux430ux43bux44bux433-ux431ux430ux442ux43bux430ux445-25-ux43cux438ux43d}{%
\subsection{Алхам 1 --- Суурь ба аргачлалыг батлах (25 мин)}\label{ux430ux43bux445ux430ux43c-1-ux441ux443ux443ux440ux44c-ux431ux430-ux430ux440ux433ux430ux447ux43bux430ux43bux44bux433-ux431ux430ux442ux43bux430ux445-25-ux43cux438ux43d}}

Загваргүйгээр хэмжилтийн аргачлалыг шалгана:

\begin{Shaded}
\begin{Highlighting}[]
\BuiltInTok{cd}\NormalTok{ \textasciitilde{}/cnc302}
\ExtensionTok{python3}\NormalTok{ lab06/benchmark\_inference.py }\AttributeTok{{-}{-}backend}\NormalTok{ synthetic }\AttributeTok{{-}{-}runs}\NormalTok{ 500 }\DataTypeTok{\textbackslash{}}
    \AttributeTok{{-}{-}threads}\NormalTok{ 1 }\AttributeTok{{-}{-}label}\NormalTok{ baseline{-}t1 }\AttributeTok{{-}{-}csv}\NormalTok{ lab06/out/bench.csv}
\end{Highlighting}
\end{Shaded}

\textbf{Халаалтын нөлөөг харна.} \texttt{-\/-warmup\ 0} ба \texttt{-\/-warmup\ 50}-ыг харьцуулна:

\begin{Shaded}
\begin{Highlighting}[]
\ExtensionTok{python3}\NormalTok{ lab06/benchmark\_inference.py }\AttributeTok{{-}{-}backend}\NormalTok{ synthetic }\AttributeTok{{-}{-}runs}\NormalTok{ 200 }\AttributeTok{{-}{-}warmup}\NormalTok{ 0 }\DataTypeTok{\textbackslash{}}
    \AttributeTok{{-}{-}label}\NormalTok{ no{-}warmup }\AttributeTok{{-}{-}csv}\NormalTok{ lab06/out/bench.csv}
\ExtensionTok{python3}\NormalTok{ lab06/benchmark\_inference.py }\AttributeTok{{-}{-}backend}\NormalTok{ synthetic }\AttributeTok{{-}{-}runs}\NormalTok{ 200 }\AttributeTok{{-}{-}warmup}\NormalTok{ 50 }\DataTypeTok{\textbackslash{}}
    \AttributeTok{{-}{-}label}\NormalTok{ warm{-}50 }\AttributeTok{{-}{-}csv}\NormalTok{ lab06/out/bench.csv}
\end{Highlighting}
\end{Shaded}

\texttt{lat\_\allowbreak{}max\_\allowbreak{}ms} ба \texttt{lat\_\allowbreak{}p99\_\allowbreak{}ms}-ийн ялгааг тэмдэглэ. \textbf{Хүснэгт 2}-ын эхний хоёр мөр.

\hypertarget{ux430ux43bux445ux430ux43c-2-float32-ux431ux430-int8-50-ux43cux438ux43d}{%
\subsection{Алхам 2 --- float32 ба int8 (50 мин)}\label{ux430ux43bux445ux430ux43c-2-float32-ux431ux430-int8-50-ux43cux438ux43d}}

\begin{Shaded}
\begin{Highlighting}[]
\ExtensionTok{python3}\NormalTok{ lab06/benchmark\_inference.py }\AttributeTok{{-}{-}backend}\NormalTok{ tflite }\DataTypeTok{\textbackslash{}}
    \AttributeTok{{-}{-}model}\NormalTok{ lab06/models/model\_float32.tflite }\AttributeTok{{-}{-}runs}\NormalTok{ 300 }\AttributeTok{{-}{-}threads}\NormalTok{ 1 }\DataTypeTok{\textbackslash{}}
    \AttributeTok{{-}{-}label}\NormalTok{ float32{-}t1 }\AttributeTok{{-}{-}csv}\NormalTok{ lab06/out/bench.csv}

\ExtensionTok{python3}\NormalTok{ lab06/benchmark\_inference.py }\AttributeTok{{-}{-}backend}\NormalTok{ tflite }\DataTypeTok{\textbackslash{}}
    \AttributeTok{{-}{-}model}\NormalTok{ lab06/models/model\_int8.tflite }\AttributeTok{{-}{-}runs}\NormalTok{ 300 }\AttributeTok{{-}{-}threads}\NormalTok{ 1 }\DataTypeTok{\textbackslash{}}
    \AttributeTok{{-}{-}label}\NormalTok{ int8{-}t1 }\AttributeTok{{-}{-}csv}\NormalTok{ lab06/out/bench.csv}
\end{Highlighting}
\end{Shaded}

\textbf{Хүснэгт 1} ба \textbf{Хүснэгт 2}-ыг бөглөнө. Тооцоолно:

\begin{verbatim}
Хурдны ашиг    = саатал(float32) / саатал(int8)
Санах ойн ашиг = хэмжээ(float32) / хэмжээ(int8)
Нарийвчлалын алдагдал = acc(float32) − acc(int8)     ← Edge Impulse Studio-гоос
\end{verbatim}

\begin{quote}
\textbf{Хэрэв int8 нь float32-оос хурдан БИШ бол} энэ нь буруу биш. Жижиг загварт квантчилалын нэмэлт зардал (де/квантчилал) ашгийг идэж магадгүй. Үүнийг тайлбарлах нь чухал ажиглалт.
\end{quote}

\hypertarget{ux430ux43bux445ux430ux43c-3-ux443ux440ux441ux433ux430ux43bux44bux43d-ux442ux43eux43e-ux431ux430-ux437ux44dux440ux44dux433ux446ux44dux44d-ux430ux436ux438ux43bux43bux430ux433ux430ux430-30-ux43cux438ux43d}{%
\subsection{Алхам 3 --- Урсгалын тоо ба зэрэгцээ ажиллагаа (30 мин)}\label{ux430ux43bux445ux430ux43c-3-ux443ux440ux441ux433ux430ux43bux44bux43d-ux442ux43eux43e-ux431ux430-ux437ux44dux440ux44dux433ux446ux44dux44d-ux430ux436ux438ux43bux43bux430ux433ux430ux430-30-ux43cux438ux43d}}

\begin{Shaded}
\begin{Highlighting}[]
\ControlFlowTok{for}\NormalTok{ T }\KeywordTok{in}\NormalTok{ 1 2 4}\KeywordTok{;} \ControlFlowTok{do}
  \ExtensionTok{python3}\NormalTok{ lab06/benchmark\_inference.py }\AttributeTok{{-}{-}backend}\NormalTok{ tflite }\DataTypeTok{\textbackslash{}}
      \AttributeTok{{-}{-}model}\NormalTok{ lab06/models/model\_int8.tflite }\AttributeTok{{-}{-}runs}\NormalTok{ 300 }\AttributeTok{{-}{-}threads} \VariableTok{$T} \DataTypeTok{\textbackslash{}}
      \AttributeTok{{-}{-}label}\NormalTok{ int8{-}t}\VariableTok{$T} \AttributeTok{{-}{-}csv}\NormalTok{ lab06/out/bench.csv}
\ControlFlowTok{done}
\end{Highlighting}
\end{Shaded}

\textbf{Хүснэгт 3}-ыг бөглөнө. Дараах асуултад хариулна:

\begin{itemize}
\tightlist
\item
  4 урсгал 1 урсгалаас 4 дахин хурдан болов уу? Хэрэв үгүй бол яагаад?
\item
  \texttt{cpu\_efficiency} (ашиглагдсан цөмийн тоо) ямар байв?
\item
  \textbf{Ирмэгийн зангилаа дээр 4 цөмийг бүгдийг AI-д өгөх нь зөв үү?} Санамж: тэр Pi дээр EMQX, ThingsBoard, InfluxDB бас ажиллаж байгааг мартуузай.
\end{itemize}

\textbf{Ачаалалтай үеийн хэмжилт.} Стек ажиллаж, 50 төхөөрөмж өгөгдөл илгээж байхад давтана:

\begin{Shaded}
\begin{Highlighting}[]
\CommentTok{\# зөөврийн компьютер дээр}
\ExtensionTok{python3}\NormalTok{ tools/sim\_device.py }\AttributeTok{{-}{-}host} \VariableTok{$H} \AttributeTok{{-}{-}devices}\NormalTok{ 50 }\AttributeTok{{-}{-}interval}\NormalTok{ 1 }\KeywordTok{\&}

\CommentTok{\# Pi дээр}
\ExtensionTok{python3}\NormalTok{ lab06/benchmark\_inference.py }\AttributeTok{{-}{-}backend}\NormalTok{ tflite }\DataTypeTok{\textbackslash{}}
    \AttributeTok{{-}{-}model}\NormalTok{ lab06/models/model\_int8.tflite }\AttributeTok{{-}{-}runs}\NormalTok{ 300 }\AttributeTok{{-}{-}threads}\NormalTok{ 1 }\DataTypeTok{\textbackslash{}}
    \AttributeTok{{-}{-}label}\NormalTok{ int8{-}t1{-}loaded }\AttributeTok{{-}{-}csv}\NormalTok{ lab06/out/bench.csv}
\end{Highlighting}
\end{Shaded}

Чөлөөт vs ачаалалтай үеийн p99-ийн ялгаа: ule{1.5cm}{0.4pt} \%

\hypertarget{ux430ux43bux445ux430ux43c-4-ux445ux443ux440ux434ux430ux441ux433ux443ux443ux440-raspberry-pi-ai-kit-40-ux43cux438ux43d}{%
\subsection{Алхам 4 --- Хурдасгуур (Raspberry Pi AI Kit) (40 мин)}\label{ux430ux43bux445ux430ux43c-4-ux445ux443ux440ux434ux430ux441ux433ux443ux443ux440-raspberry-pi-ai-kit-40-ux43cux438ux43d}}

\begin{quote}
AI Kit тэнхимийн хамтын хэрэглээ тул ээлжээр ажиллана. Хүлээх зуураа Алхам 5-ыг эхлүүлж болно.
\end{quote}

\textbf{4.1 Тоног төхөөрөмжийг шалгах:}

\begin{Shaded}
\begin{Highlighting}[]
\ExtensionTok{lspci} \KeywordTok{|} \FunctionTok{grep} \AttributeTok{{-}i}\NormalTok{ hailo}
\ExtensionTok{hailortcli}\NormalTok{ fw{-}control identify}
\end{Highlighting}
\end{Shaded}

\textbf{4.2 Edge Impulse-ийн Linux runner-ээр:}

\begin{Shaded}
\begin{Highlighting}[]
\ExtensionTok{python3}\NormalTok{ lab06/benchmark\_inference.py }\AttributeTok{{-}{-}backend}\NormalTok{ eim }\DataTypeTok{\textbackslash{}}
    \AttributeTok{{-}{-}model}\NormalTok{ lab06/models/anomaly.eim }\AttributeTok{{-}{-}runs}\NormalTok{ 300 }\DataTypeTok{\textbackslash{}}
    \AttributeTok{{-}{-}label}\NormalTok{ eim{-}cpu }\AttributeTok{{-}{-}csv}\NormalTok{ lab06/out/bench.csv}
\end{Highlighting}
\end{Shaded}

\textbf{4.3} Хэрэв таны Edge Impulse төсөл Hailo зорилтыг дэмждэг бол хурдасгасан хувилбарыг татаж дахин хэмжинэ (\texttt{-\/-label\ eim-hailo}).

\begin{quote}
\textbf{Дэмжихгүй бол:} Энэ нь бүрэн хүлээн зөвшөөрөгдөх үр дүн. Тайландаа \textbf{яагаад дэмжихгүй байгааг} судалж бичнэ: загварын архитектур Hailo-гийн үйлдлийн жагсаалтад нийцэхгүй байж болно, эсвэл сургалтын хэлбэр (K-means anomaly) нь NPU-д хөрвүүлэгддэггүй. Энэ бол ирмэгийн AI-ийн бодит хязгаарлалт --- \textbf{хурдасгуур бүх загварыг хурдасгадаггүй}.
\end{quote}

\textbf{Хүснэгт 4}-ийг бөглөнө.

\hypertarget{ux430ux43bux445ux430ux43c-5-ux43fux43bux430ux442ux444ux43eux440ux43cux434-ux43dux44dux433ux442ux433ux44dux445-40-ux43cux438ux43d}{%
\subsection{Алхам 5 --- Платформд нэгтгэх (40 мин)}\label{ux430ux43bux445ux430ux43c-5-ux43fux43bux430ux442ux444ux43eux440ux43cux434-ux43dux44dux433ux442ux433ux44dux445-40-ux43cux438ux43d}}

Дүгнэлтийг зөвхөн хэмжээд орхихгүй, \textbf{системд буцаан холбоно}.

\begin{Shaded}
\begin{Highlighting}[]
\CommentTok{\# терминал 1 (Pi) — ирмэгийн дүгнэлт}
\ExtensionTok{python3}\NormalTok{ lab06/anomaly\_publish.py }\AttributeTok{{-}{-}host}\NormalTok{ localhost }\AttributeTok{{-}{-}backend}\NormalTok{ tflite }\DataTypeTok{\textbackslash{}}
    \AttributeTok{{-}{-}model}\NormalTok{ lab06/models/model\_int8.tflite }\AttributeTok{{-}{-}window}\NormalTok{ 125 }\AttributeTok{{-}{-}threshold}\NormalTok{ 0.75}

\CommentTok{\# терминал 2 (зөөврийн компьютер) — гажилтай өгөгдөл}
\ExtensionTok{python3}\NormalTok{ tools/sim\_device.py }\AttributeTok{{-}{-}host} \VariableTok{$H} \AttributeTok{{-}{-}devices}\NormalTok{ 5 }\AttributeTok{{-}{-}interval}\NormalTok{ 0.2 }\AttributeTok{{-}{-}anomaly{-}rate}\NormalTok{ 0.1}

\CommentTok{\# терминал 3 — event сувгийг сонсоно}
\ExtensionTok{mosquitto\_sub} \AttributeTok{{-}h} \VariableTok{$H} \AttributeTok{{-}t} \StringTok{\textquotesingle{}cnc302/+/+/+/+/event\textquotesingle{}} \AttributeTok{{-}v}
\end{Highlighting}
\end{Shaded}

\textbf{5.1} Лаб 4-т үүсгэсэн EMQX Rule Engine дүрмийг өргөтгөж, \texttt{event} сувгийн \texttt{type=anomaly} мессежийг \texttt{cnc302/alerts/ai} руу чиглүүлнэ.

\textbf{5.2} Лаб 5-ын Grafana самбарт ирмэгийн дохиоллын панель нэмнэ.

\textbf{5.3 Хамгийн чухал хэмжилт --- зурвасын өргөний хэмнэлт:}

\begingroup\scriptsize\begin{longtable}[]{@{}lll@{}}
\toprule
& Түүхий өгөгдөл үүлэн рүү & Зөвхөн дүгнэлтийн үр дүн \\
\midrule
\endhead
\bottomrule
\endlastfoot
Мессеж/секунд & 5 төх × 5 мсж/с = 25 & зөвхөн гажил үед ≈ 0.5 \\
Байт/мессеж & \textasciitilde140 & \textasciitilde180 \\
\textbf{Байт/секунд} & & \\
\textbf{Өдөрт (MiB)} & & \\
\textbf{Хэмнэлт \%} & & \\
\end{longtable}\endgroup

\hypertarget{ux430ux43bux445ux430ux43c-6-ux448ux438ux43dux436ux438ux43bux433ux44dux44d-ux431ux430-ux448ux438ux439ux434ux432ux44dux440-25-ux43cux438ux43d}{%
\subsection{Алхам 6 --- Шинжилгээ ба шийдвэр (25 мин)}\label{ux430ux43bux445ux430ux43c-6-ux448ux438ux43dux436ux438ux43bux433ux44dux44d-ux431ux430-ux448ux438ux439ux434ux432ux44dux440-25-ux43cux438ux43d}}

\begin{Shaded}
\begin{Highlighting}[]
\ExtensionTok{python3} \AttributeTok{{-}} \OperatorTok{\textless{}\textless{}\textquotesingle{}EOF\textquotesingle{}}
\StringTok{import csv}
\StringTok{rows=list(csv.DictReader(open(\textquotesingle{}lab06/out/bench.csv\textquotesingle{})))}
\StringTok{w=max(len(r[\textquotesingle{}label\textquotesingle{}]) for r in rows)}
\StringTok{print(f"\{\textquotesingle{}label\textquotesingle{}:\textless{}\{w\}\} \{\textquotesingle{}p50\textquotesingle{}:\textgreater{}8\} \{\textquotesingle{}p99\textquotesingle{}:\textgreater{}8\} \{\textquotesingle{}MiB\textquotesingle{}:\textgreater{}7\} \{\textquotesingle{}KiB\textquotesingle{}:\textgreater{}8\} \{\textquotesingle{}инф/с\textquotesingle{}:\textgreater{}9\}")}
\StringTok{for r in rows:}
\StringTok{    print(f"\{r[\textquotesingle{}label\textquotesingle{}]:\textless{}\{w\}\} \{r[\textquotesingle{}lat\_p50\_ms\textquotesingle{}]:\textgreater{}8\} \{r[\textquotesingle{}lat\_p99\_ms\textquotesingle{}]:\textgreater{}8\} "}
\StringTok{          f"\{r[\textquotesingle{}mem\_model\_mib\textquotesingle{}]:\textgreater{}7\} \{r[\textquotesingle{}model\_kib\textquotesingle{}]:\textgreater{}8\} \{r[\textquotesingle{}throughput\_infer\_s\textquotesingle{}]:\textgreater{}9\}")}
\OperatorTok{EOF}
\end{Highlighting}
\end{Shaded}

\textbf{Хүснэгт 5} (байршуулалтын шийдвэр)-ийг бөглөж commit хийнэ:

\begin{Shaded}
\begin{Highlighting}[]
\FunctionTok{git}\NormalTok{ add }\AttributeTok{{-}A} \KeywordTok{\&\&} \FunctionTok{git}\NormalTok{ status      }\CommentTok{\# models/, data/, out/ орохгүй!}
\FunctionTok{git}\NormalTok{ commit }\AttributeTok{{-}m} \StringTok{"Лаб 6: ирмэгийн AI{-}ийн гүйцэтгэлийн хэмжилт"}
\FunctionTok{git}\NormalTok{ tag lab06{-}done }\KeywordTok{\&\&} \FunctionTok{git}\NormalTok{ push origin main }\AttributeTok{{-}{-}tags}
\end{Highlighting}
\end{Shaded}

\hypertarget{ux445ux44dux43cux436ux438ux43bux442ux438ux439ux43d-ux445ux4afux441ux43dux44dux433ux442ux4afux4afux434}{%
\section{5. Хэмжилтийн хүснэгтүүд}\label{ux445ux44dux43cux436ux438ux43bux442ux438ux439ux43d-ux445ux4afux441ux43dux44dux433ux442ux4afux4afux434}}

\hypertarget{ux445ux4afux441ux43dux44dux433ux442-1-ux437ux430ux433ux432ux430ux440ux44bux43d-ux43cux44dux434ux44dux44dux43bux44dux43b-edge-impulse-studio-ux433ux43eux43eux441}{%
\subsection{Хүснэгт 1 --- Загварын мэдээлэл (Edge Impulse Studio-гоос)}\label{ux445ux4afux441ux43dux44dux433ux442-1-ux437ux430ux433ux432ux430ux440ux44bux43d-ux43cux44dux434ux44dux44dux43bux44dux43b-edge-impulse-studio-ux433ux43eux43eux441}}

\begingroup\scriptsize\begin{longtable}[]{@{}lll@{}}
\toprule
Үзүүлэлт & float32 & int8 \\
\midrule
\endhead
\bottomrule
\endlastfoot
Загварын төрөл & & \\
Оролтын цонх (мс / дээж) & & \\
Нарийвчлал (Model testing) & & \\
Studio-гийн таамагласан саатал & & \\
Studio-гийн таамагласан RAM & & \\
Файлын хэмжээ (KiB) & & \\
\end{longtable}\endgroup

\hypertarget{ux445ux4afux441ux43dux44dux433ux442-2-ux445ux44dux43cux436ux441ux44dux43d-ux433ux4afux439ux446ux44dux442ux433ux44dux43b-1-ux443ux440ux441ux433ux430ux43b}{%
\subsection{Хүснэгт 2 --- Хэмжсэн гүйцэтгэл (1 урсгал)}\label{ux445ux4afux441ux43dux44dux433ux442-2-ux445ux44dux43cux436ux441ux44dux43d-ux433ux4afux439ux446ux44dux442ux433ux44dux43b-1-ux443ux440ux441ux433ux430ux43b}}

\begingroup\scriptsize\begin{longtable}[]{@{}
  >{\raggedright\arraybackslash}p{(\columnwidth - 14\tabcolsep) * \real{0.1250}}
  >{\raggedright\arraybackslash}p{(\columnwidth - 14\tabcolsep) * \real{0.1250}}
  >{\raggedright\arraybackslash}p{(\columnwidth - 14\tabcolsep) * \real{0.1250}}
  >{\raggedright\arraybackslash}p{(\columnwidth - 14\tabcolsep) * \real{0.1250}}
  >{\raggedright\arraybackslash}p{(\columnwidth - 14\tabcolsep) * \real{0.1250}}
  >{\raggedright\arraybackslash}p{(\columnwidth - 14\tabcolsep) * \real{0.1250}}
  >{\raggedright\arraybackslash}p{(\columnwidth - 14\tabcolsep) * \real{0.1250}}
  >{\raggedright\arraybackslash}p{(\columnwidth - 14\tabcolsep) * \real{0.1250}}@{}}
\toprule
\begin{minipage}[b]{\linewidth}\raggedright
Шошго
\end{minipage} & \begin{minipage}[b]{\linewidth}\raggedright
p50 (мс)
\end{minipage} & \begin{minipage}[b]{\linewidth}\raggedright
p95 (мс)
\end{minipage} & \begin{minipage}[b]{\linewidth}\raggedright
p99 (мс)
\end{minipage} & \begin{minipage}[b]{\linewidth}\raggedright
max (мс)
\end{minipage} & \begin{minipage}[b]{\linewidth}\raggedright
σ (мс)
\end{minipage} & \begin{minipage}[b]{\linewidth}\raggedright
Санах ой (MiB)
\end{minipage} & \begin{minipage}[b]{\linewidth}\raggedright
Чадвар (инф/с)
\end{minipage} \\
\midrule
\endhead
\bottomrule
\endlastfoot
baseline-t1 (synthetic) & & & & & & & \\
no-warmup & & & & & & & \\
warm-50 & & & & & & & \\
float32-t1 & & & & & & & \\
int8-t1 & & & & & & & \\
int8-t1-loaded & & & & & & & \\
\end{longtable}\endgroup

Квантчилалын хурдны ашиг: ule{1.5cm}{0.4pt} × · Санах ойн ашиг: ule{1.5cm}{0.4pt} × · Нарийвчлалын алдагдал: ule{1.5cm}{0.4pt} \%

\hypertarget{ux445ux4afux441ux43dux44dux433ux442-3-ux443ux440ux441ux433ux430ux43bux44bux43d-ux442ux43eux43e}{%
\subsection{Хүснэгт 3 --- Урсгалын тоо}\label{ux445ux4afux441ux43dux44dux433ux442-3-ux443ux440ux441ux433ux430ux43bux44bux43d-ux442ux43eux43e}}

\begingroup\scriptsize\begin{longtable}[]{@{}
  >{\raggedright\arraybackslash}p{(\columnwidth - 10\tabcolsep) * \real{0.1667}}
  >{\raggedright\arraybackslash}p{(\columnwidth - 10\tabcolsep) * \real{0.1667}}
  >{\raggedright\arraybackslash}p{(\columnwidth - 10\tabcolsep) * \real{0.1667}}
  >{\raggedright\arraybackslash}p{(\columnwidth - 10\tabcolsep) * \real{0.1667}}
  >{\raggedright\arraybackslash}p{(\columnwidth - 10\tabcolsep) * \real{0.1667}}
  >{\raggedright\arraybackslash}p{(\columnwidth - 10\tabcolsep) * \real{0.1667}}@{}}
\toprule
\begin{minipage}[b]{\linewidth}\raggedright
Урсгал
\end{minipage} & \begin{minipage}[b]{\linewidth}\raggedright
p50 (мс)
\end{minipage} & \begin{minipage}[b]{\linewidth}\raggedright
p99 (мс)
\end{minipage} & \begin{minipage}[b]{\linewidth}\raggedright
Чадвар (инф/с)
\end{minipage} & \begin{minipage}[b]{\linewidth}\raggedright
CPU ашиглалт (цөм)
\end{minipage} & \begin{minipage}[b]{\linewidth}\raggedright
Хурдны ашиг vs t1
\end{minipage} \\
\midrule
\endhead
\bottomrule
\endlastfoot
1 & & & & & 1.00× \\
2 & & & & & \\
4 & & & & & \\
\end{longtable}\endgroup

4 урсгал 4× хурдан болов уу: ule{1.5cm}{0.4pt} Шалтгаан: ule{1.5cm}{0.4pt}

\hypertarget{ux445ux4afux441ux43dux44dux433ux442-4-ux445ux443ux440ux434ux430ux441ux433ux443ux443ux440ux44bux43d-ux445ux430ux440ux44cux446ux443ux443ux43bux430ux43bux442}{%
\subsection{Хүснэгт 4 --- Хурдасгуурын харьцуулалт}\label{ux445ux4afux441ux43dux44dux433ux442-4-ux445ux443ux440ux434ux430ux441ux433ux443ux443ux440ux44bux43d-ux445ux430ux440ux44cux446ux443ux443ux43bux430ux43bux442}}

\begingroup\scriptsize\begin{longtable}[]{@{}lllll@{}}
\toprule
Гүйцэтгэх орчин & p50 (мс) & p99 (мс) & Санах ой (MiB) & Тэмдэглэл \\
\midrule
\endhead
\bottomrule
\endlastfoot
CPU, float32 & & & & \\
CPU, int8 & & & & \\
Edge Impulse EIM (CPU) & & & & \\
Hailo-8L (AI Kit) & & & & дэмжигдсэн үү: \\
\end{longtable}\endgroup

\hypertarget{ux445ux4afux441ux43dux44dux433ux442-5-ux431ux430ux439ux440ux448ux443ux443ux43bux430ux43bux442ux44bux43d-ux448ux438ux439ux434ux432ux44dux440}{%
\subsection{Хүснэгт 5 --- Байршуулалтын шийдвэр}\label{ux445ux4afux441ux43dux44dux433ux442-5-ux431ux430ux439ux440ux448ux443ux443ux43bux430ux43bux442ux44bux43d-ux448ux438ux439ux434ux432ux44dux440}}

\begingroup\scriptsize\begin{longtable}[]{@{}lll@{}}
\toprule
Асуулт & Хариулт & Тоон үндэслэл \\
\midrule
\endhead
\bottomrule
\endlastfoot
Сонгосон загвар ба формат & & \\
Урсгалын тоо & & \\
100 Гц мэдрэгчид нийцэх үү & & p99 = \_\_\_ мс \\
Хэдэн төхөөрөмжийн урсгалыг нэг Pi дүгнэж чадах вэ & & \\
Ирмэг эсвэл үүл --- сонголт ба шалтгаан & & \\
Зурвасын өргөний хэмнэлт & & \_\_\_ \% \\
Загварыг хэрхэн шинэчлэх вэ & & Лаб 2-ын OTA-тай хэрхэн холбогдох \\
\end{longtable}\endgroup

\hypertarget{ux445ux44fux43dux430ux43bux442ux44bux43d-ux430ux441ux443ux443ux43bux442}{%
\section{6. Хяналтын асуулт}\label{ux445ux44fux43dux430ux43bux442ux44bux43d-ux430ux441ux443ux443ux43bux442}}

\begin{enumerate}
\def\labelenumi{\arabic{enumi}.}
\tightlist
\item
  Халаалтгүй ба халаалттай хэмжилтийн \texttt{max} утга хэр ялгаатай байв? Яагаад анхны дүгнэлтүүд удаан вэ?
\item
  Яагаад p50 биш p99-ээр төлөвлөх ёстой вэ? p99 хэтэрвэл бодит системд юу болох вэ?
\item
  int8 квантчилал таны загварт хэр их ашиг өгөв? Хэрэв бага бол яагаад?
\item
  Ачаалалтай үед саатал өссөн үү? Энэ нь ирмэгийн зангилааны төлөвлөлтөд ямар утгатай вэ?
\item
  Дүгнэлтийг ирмэг дээр хийснээр хэдэн хувийн зурвасын өргөн хэмнэгдэв? 1000 төхөөрөмж дээр өдөрт хэдэн GiB вэ?
\item
  Ирмэг дээрх загварыг хэрхэн шинэчлэх вэ? Лаб 2-т үзсэн OTA механизм загварт тохирох уу? Ямар онцлог шаардлага гарах вэ?
\item
  Загварын нарийвчлал байршуулсны дараа буурч болох уу? (Санамж: \textbf{model drift}.) Үүнийг хэрхэн илрүүлэх вэ? Лаб 5-ын Grafana дохиолол энд яаж тусалж чадах вэ?
\item
  Hailo хурдасгуур таны загварыг дэмжсэн үү? Дэмжээгүй бол шалтгааныг тайлбарла.
\item
  Холбооны сургалт (federated learning) энэ хэрэглээнд ямар давуу тал өгөх байсан бэ?
\end{enumerate}

\hypertarget{ux442ux430ux439ux43bux430ux43dux433ux438ux439ux43d-ux448ux430ux430ux440ux434ux43bux430ux433ux430}{%
\section{7. Тайлангийн шаардлага}\label{ux442ux430ux439ux43bux430ux43dux433ux438ux439ux43d-ux448ux430ux430ux440ux434ux43bux430ux433ux430}}

\begin{itemize}
\tightlist
\item[$\square$]
  Хүснэгт 1--5 бүрэн
\item[$\square$]
  \texttt{lab06/out/bench.csv} (бүх хэмжилт нэг файлд)
\item[$\square$]
  Edge Impulse төслийн холбоос эсвэл дэлгэцийн зураг (Model testing хуудас)
\item[$\square$]
  Ирмэгийн дохиолол Grafana дээр харагдаж буй нотолгоо
\item[$\square$]
  Зурвасын өргөний хэмнэлтийн тооцоо
\item[$\square$]
  Хяналтын 9 асуултын хариулт
\item[$\square$]
  Git tag \texttt{lab06-done}
\end{itemize}

\hypertarget{ux43eux43dux43eux448ux438ux43bux433ux43eux43e}{%
\section{8. Оношилгоо}\label{ux43eux43dux43eux448ux438ux43bux433ux43eux43e}}

\begingroup\scriptsize\begin{longtable}[]{@{}
  >{\raggedright\arraybackslash}p{(\columnwidth - 4\tabcolsep) * \real{0.3333}}
  >{\raggedright\arraybackslash}p{(\columnwidth - 4\tabcolsep) * \real{0.3333}}
  >{\raggedright\arraybackslash}p{(\columnwidth - 4\tabcolsep) * \real{0.3333}}@{}}
\toprule
\begin{minipage}[b]{\linewidth}\raggedright
Шинж тэмдэг
\end{minipage} & \begin{minipage}[b]{\linewidth}\raggedright
Шалтгаан
\end{minipage} & \begin{minipage}[b]{\linewidth}\raggedright
Шийдэл
\end{minipage} \\
\midrule
\endhead
\bottomrule
\endlastfoot
\texttt{TFLite\ орчин\ олдсонгүй} & сан суулгаагүй & \texttt{pip3\ install\ -\/-break-system-packages\ ai-edge-litert} \\
\texttt{.eim} ажиллахгүй & гүйцэтгэх эрхгүй эсвэл буруу архитектур & \texttt{chmod\ +x}; AARCH64 хувилбар татсан эсэхээ шалга \\
int8 float32-оос удаан & загвар хэт жижиг, де/квантчилалын зардал давамгайлав & хэвийн үзэгдэл --- тайлбарла \\
Саатал маш хэлбэлзэлтэй (σ өндөр) & Pi дээр өөр ачаалал байна & \texttt{htop}; стекийг түр зогсоож давт \\
\texttt{lspci} Hailo харуулахгүй & AI Kit холбогдоогүй эсвэл PCIe идэвхгүй & \texttt{sudo\ raspi-config} → PCIe Gen 3 \\
Санах ойн хэмжилт 0 & \texttt{/proc/self/status} уншигдахгүй & контейнер дотор биш, Pi дээр шууд ажиллуул \\
\texttt{anomaly\_publish.py} дүгнэлт хийхгүй & цонх дүүрэхгүй байна & \texttt{-\/-window}-ыг багасга эсвэл \texttt{-\/-interval}-ыг хурдасга \\
\end{longtable}\endgroup

\hypertarget{ux4afux43dux44dux43bux433ux44dux44dux43dux438ux439-ux448ux430ux43bux433ux443ux443ux440-10-ux43eux43dux43eux43e}{%
\section{9. Үнэлгээний шалгуур (10 оноо)}\label{ux4afux43dux44dux43bux433ux44dux44dux43dux438ux439-ux448ux430ux43bux433ux443ux443ux440-10-ux43eux43dux43eux43e}}

\begingroup\scriptsize\begin{longtable}[]{@{}
  >{\raggedright\arraybackslash}p{(\columnwidth - 2\tabcolsep) * \real{0.5000}}
  >{\raggedright\arraybackslash}p{(\columnwidth - 2\tabcolsep) * \real{0.5000}}@{}}
\toprule
\begin{minipage}[b]{\linewidth}\raggedright
Шалгуур
\end{minipage} & \begin{minipage}[b]{\linewidth}\raggedright
Оноо
\end{minipage} \\
\midrule
\endhead
\bottomrule
\endlastfoot
Загвар Pi дээр байршиж ажиллаж байна & 2 \\
float32/int8, урсгалын тоо, ачаалалтай/чөлөөт харьцуулалт хэмжигдсэн & 3 \\
Хурдасгуурын туршилт (дэмжигдээгүй бол шалтгаан судлагдсан) & 1 \\
Дүгнэлт платформд нэгтгэгдэж, дохиолол ажиллаж байна & 2 \\
Зурвасын өргөний тооцоо ба байршуулалтын шийдвэр & 2 \\
\end{longtable}\endgroup

\textbf{Загваргүйгээр (\texttt{-\/-backend\ synthetic}) гүйцэтгэсэн бол} дээд тал нь 7 оноо.
